\subsection{Recovering the equilibrium path}\label{\refsec:PEEpath}
Algorithms \ref{\refalg:LOG:gridsearch} and \ref{\refalg:CRRA:grid} identify policy \emph{functions}, not a path: at each $t$ they return $\tau_t$ as a function of the state inherited from $t-1$. To report an actual politico-economic equilibrium we still have to pin down the state entering the first period and then walk the functions forward. This step is unavoidable in both preference cases here, since neither is free of endogenous states.

\smalltitle{The initial state is a scalar fixed point.}
Following \cref{\refsec:EE}, we take the pre-determined period as given by a steady state assumption: the generation that is old in the first active period behaved as if it had always faced the policy it meets there. That policy is itself what we are solving for, which is what makes this a fixed point rather than a substitution.

Write the first active period as $t=1$ and let $\tau_1$ denote the tax chosen there. Under the steady state assumption, a constant policy $\tau_1$ delivers a steady-state coefficient function $\Theta_s(\tau_1)$ from \eqref{\refeq:steadystate_LOG} or \eqref{\refeq:steadystate_CRRA} as appropriate, and hence
\begin{subequations}\label{\refeq:initialState}
\begin{align}
    \iota_0 &= \iota_0\left(B_1^0(\tau_1),\, \Theta_s(\tau_1),\, \tau_1\right), \label{\refeq:initialState:iota}\\
    s_0 &= s^*(\tau_1), \label{\refeq:initialState:s}
\end{align}
\end{subequations}
where \eqref{\refeq:initialState:iota} is \eqref{\refeq:auxiliary:s0_s} evaluated at a constant policy and \eqref{\refeq:initialState:s} is the steady-state savings level. Both are functions of $\tau_1$ alone. Closing the loop with the solved first-period policy function gives the residual
\begin{align}\label{\refeq:initialFixedPoint}
    \tau_1 - \tau^{1}\left(\iota_0(\tau_1)\right) = 0 \quad\text{(log)}, \qquad\qquad
    \tau_1 - \tau^{1}\left(s_0(\tau_1),\, \iota_0(\tau_1)\right) = 0 \quad\text{(CRRA)},
\end{align}
a scalar problem in both cases. We search over $\tau_1$ rather than over the states themselves, and deliberately so: $\tau_1$ has clean, known bounds $[l,u]$, whereas any bracket for $s_0$ or $\iota_0$ would have to be inferred from the same steady-state map we are trying to invert. Note also that under log preferences $s_0$ plays no part in \eqref{\refeq:initialFixedPoint}: it is not a political state there, so it is resolved after the fact from \eqref{\refeq:initialState:s} at the converged $\tau_1$ and used only to report levels.

\smalltitle{The bracket is guaranteed, and that is the problem.}
Evaluating $\tau^1(\cdot)$ in \eqref{\refeq:initialFixedPoint} means clipping it back into $[l,u]$, since a policy function returns a tax. That makes $[l,u]$ a bracket for free: the residual is weakly negative at $l$ and weakly positive at $u$ whatever the policy function does. It is tempting to stop there and hand the interval to a bracketed scalar solver. That is wrong, and the reason generalises past this particular fixed point.

The map $\tau_1\mapsto(s_0,\iota_0)$ of \eqref{\refeq:initialState} does not stay inside the grids the policy functions were solved on. As $\tau_1\rightarrow 1$ the steady state degenerates, and $\iota_0$ diverges --- not because informal savings explode but because their denominator, formal savings, collapses --- reaching order $10^3$ in our calibration against an upper grid bound $u_{\iota}=2$. Over that region $\tau^1(\cdot)$ is being extrapolated, and whatever it returns is then clipped. If the extrapolated value exceeds $u$, the residual at $\tau_1 = u$ is $u-u$: \emph{exactly} zero. The clip has manufactured a root at the bracket end, a bracketed solver accepts it, and the reported equilibrium is the degenerate steady state, with the interior never evaluated. Which way this falls is not a property of the model. In our calibration the extrapolated policy comes back at $8.3$ when $\rho=1.15$ and clips to $u$, giving the spurious root; at $\rho=1$ it comes back negative, clips to $l$, and leaves a residual of about $1$ with no spurious root at all. A clip that manufactures a bracket can equally manufacture a root, and every bounded policy function in this paper is clipped the same way.

Two corrections follow, neither specific to the initial period. First, the admissible set of $\tau_1$ is not $[l,u]$ but the subset of it on which $(s_0,\iota_0)$ lies inside $\mathcal{S}\times\mathcal{S}_0$ --- in the log case, on which $\iota_0\in\mathcal{S}_0$. Outside that subset the residual should be left undefined rather than evaluated, on exactly the principle applied to $\iota_t$ throughout \cref{\refsec:PEELOG}: a state that leaves the grid is reported, never clipped back onto it. Second, \eqref{\refeq:initialFixedPoint} need not have a unique root, and nothing in the problem says it does. We therefore evaluate the residual on $\mathcal{T}$, discard the nodes whose implied state leaves the grids, locate every sign change among those that survive, and refine only the lowest --- the convention already used for the multiple roots of \eqref{\refeq:stateResidualLOG}, with the number of crossings reported alongside the solution. If no sign change survives, either the state grids are too narrow for the steady states $\mathcal{T}$ implies, or no constant policy reproduces itself through the first period's policy function; both deserve to fail loudly rather than be papered over. In our calibration the masked residual is monotone with a single crossing, at $\tau_1 = 0.145$ under log preferences and $\tau_1 = 0.153$ at $\rho = 1.15$.

\smalltitle{Forward simulation.}
With $(s_0,\iota_0)$ in hand, walk forward for $t = 1,...,T$, reading the tax off the solved policy function and the next state off the solved state policy functions:
\begin{align}\label{\refeq:forwardSim}
    \begin{array}{lll}
    \text{(log)} & \tau_t = \tau^t\left(\iota_{t-1}\right), & \iota_t = \iota_t\left(\tau_t\right), \\[4pt]
    \text{(CRRA)} & \tau_t = \tau^t\left(s_{t-1},\iota_{t-1}\right), \qquad & s_t = s_t\left(\tau_t, s_{t-1}\right),\quad \iota_t = \iota_t\left(\tau_t,s_t\right).
    \end{array}
\end{align}
\eqref{\refeq:forwardSim} is written with the state transitions as functions of the tax --- $\iota_t(\tau_t)$ in the log case, $s_t(\tau_t,s_{t-1})$ and then $\iota_t(\tau_t,s_t)$ in the CRRA case --- and that is meant literally: they are \emph{re-solved} at the simulated $\tau_t$, not read off an interpolant of the selected states over the predetermined ones. Step 4 of \cref{\refalg:LOG:gridsearch} and \cref{\refalg:CRRA:grid} can record such an interpolant at no extra cost, and walking it is the cheaper alternative, but it interpolates a composition that the re-solve evaluates directly. What it buys is one scalar root per \emph{period}, against the one per node the recursion has already paid --- a rounding error in the cost of the whole solve. What it costs is roughly two orders of magnitude: measured against the exact re-solve described below, the re-solved walk leaves a discrepancy of order $10^{-8}$ in the log case and $10^{-7}$ in the CRRA case, against $5\cdot 10^{-6}$ and $10^{-5}$ respectively for the interpolated one. In the log case there is a second reason to prefer it, beyond accuracy: $\iota_t$ there is a function of $\tau_t$ alone, so re-solving preserves the result of \cref{\refsec:PEELOG} that the predetermined state does not enter the fixed point --- which interpolating $\iota_t$ over $\iota_{t-1}$ quietly gives up.

Three checks belong here rather than later. First, the simulated states must stay where the policy functions mean something: in the log case that is $\iota_{t-1}\in[l_{\iota},u_{\iota}]$, and in the CRRA case the stronger requirement that $(s_{t-1},\iota_{t-1})$ lie inside the reachable set $\mathcal{P}_{t-1}$ of \eqref{\refeq:reachable}. It is the state \emph{entering} each period that is tested, since that is the argument the policy function is evaluated at; the state entering the first period comes from \eqref{\refeq:initialState} rather than from any transition, so there is no reachable set to test it against. A path that leaves either is being driven by extrapolated policy and should be reported as unreliable, not returned.

The reachable-set requirement needs a word on what is actually testable. $\mathcal{P}_t$ is recorded as the image of the state grid --- a finite set of pairs, not a region --- so exact membership is not something the simulation can check. What it can check is the bounding box of that image, and the asymmetry is what makes the check worth having: the box \emph{contains} $\mathcal{P}_t$, so a simulated pair falling outside the box lies outside $\mathcal{P}_t$ for certain, while one falling inside is merely not ruled out. A containment test should be conservative in that direction and not the other.

Second, each $\tau_t$ should be checked against the corners $l$ and $u$: a path that pins to a corner for several periods usually indicates a grid or calibration problem rather than an economic result. Third, and in the CRRA case only, the simulation should test that what it reads is defined at all, not merely well positioned. A path there can become undefined without leaving $\mathcal{S}\times\mathcal{S}_0$: the policy surfaces are undefined at their infeasible nodes, and a two-dimensional interpolant cannot be restricted to its feasible nodes the way the one-dimensional interpolants of \cref{\refsec:PEELOG} can, so any cell with an infeasible corner returns an undefined value for an interior argument. The log case has no counterpart to this, which is why it is easy to omit after implementing that case first.

\smalltitle{The simulated path is never reported directly.}
The tax path from \eqref{\refeq:forwardSim} is accurate only to the resolution of the grids and to the quality of the interpolants; the allocation implied by walking the state policy functions forward inherits both errors and satisfies the equilibrium conditions only approximately. We therefore take from the simulation only $\bm{\tau}$ and the initial state $s_0$, and re-solve the economic equilibrium of \cref{\refsec:EE} \emph{exactly} at that policy path --- closed form in the log case, and the square nonlinear system in the CRRA case, warm-started at the simulated $\left(\Gamma_{s,t},h_t,s_t\right)$. Everything reported downstream, including $\iota_t$ and all consumption levels, is then read off that exact solution. The division of labour is deliberate: the grid search locates the correct branch of a problem with corners and possible multiplicity, and an exact step then recovers machine precision on a problem that has neither.

This also supplies a useful diagnostic. The $\iota_t$ path from the exact re-solve should agree with the one simulated through \eqref{\refeq:forwardSim}; a visible gap means the state grid $\mathcal{S}_0$ (or $\mathcal{S}$) is too coarse for the dynamics, and is a more informative signal than the first order condition residual, which is evaluated on the grid and therefore cannot see this error. The natural unit for ``visible'' is the spacing of the state grid itself rather than an absolute tolerance, since what the comparison asks is whether the walk resolves the state to well inside a cell. In our calibration it does, with room to spare: the simulated and re-solved paths differ by around $10^{-6}$ of a cell of $\mathcal{S}_0$ under log preferences and $10^{-4}$ of a cell under CRRA.
